Ze werd eerdaags achttien.
Om zichzelf te ontzien

van een invitatie,
koos de bange Emmy

voortijdig te vluchten.
Wat had ze te duchten

als zij naar het kasteel,
en misschien zijn bordeel,

dito werd ontboden?
De landheers ijlboden

namen moeder en zus
zonder enig excuus

mee naar zijn landerij.
Een droom daarover zei

namens een witte wief -
net als in elke brief

die ze van haar zus kreeg -
“die weg zuigt je hart leeg”.

Ze zou gaan verliezen
haar kans om te kiezen

voor een heilig leven.
Emmy’s uitvluchtstreven

werd door haar pa gesteund.
Hij had zwaar misgekleund

tijdens de sommatie
en kon via Emmy

zijn gezag herstellen.
Hij kon dan vertellen

dat hij zich deze keer
wel tegen de landheer

met verve had verzet.
De vlucht werd ingezet

vlak voor het denkbare.
Ze gaf haar dierbare

een laatste omhelzing
voordat ze op pad ging.

Haar angst bleek een uitvlucht.
Ze geloofde de klucht,

die haar brein uit kiende.
Een alles ontziende

beschermingsconstructie.
Het verhaal van Emmy

zorgde voor afscheiding.
Vervreemd van de schepping

ontstond er wantrouwen
gekloofd uit vertrouwen.

Een probleem was gesticht
en angst gaf tegenwicht.

Een zelf was gekomen.
Van jezelf afkomen

was helaas geen optie
voor de blinde Emmy.

Zonder innervaren
kon ze slechts ervaren

het enerverende
als absorberende.

Leren door emotie
te zien als reflectie

op de uiteenzetting,
die als een verklaring

dient voor het niet-weten.
Ze was God vergeten.

Dit reflectief leren
verlokte de beren

op haar ontwikkelpad,
de wikkels die ze had

nog meer aan te trekken
en niet te vertrekken.

Haar reflectief leren
werd zo rumineren.

Als ze concludeerde
dat ze rumineerde

had Emmy wel geleerd
dat veel denken haar deert.

De eerste Pinksterdag
startte met klokslag.

Haar vader aanwaaiend
en Emmy uitzwaaiend

zichtbaar bij de dorpskerk,
verried zo haar oogmerk.

Hij zocht sterk naar respect,
maar had aldus gelekt

aan valse spionnen.
Een van die pionnen

lichtte de landheer in.
Die gaf in zijn waanzin

de opdracht tot beklag
bij het bisdomsgezag.

Emmy werd vogelvrij
verklaard, omdat juist hij

door adellijke invloed
en financieel tegoed

een wet liet invoeren
dat gevluchte boeren

buiten de wet stelde
en zelfwerkend telde.

Dat die ontmoediging
automatisch inging

kon hij niet voorkomen.
Ze werd opgenomen

door een oom en tante.
Het waren galante,

maar vrezende mensen.
Ze voelden hun grenzen

als het ging om de kerk,
hun voornaamste netwerk,

boven zelfs familie.
Ze ontvingen Emmy

gedienstig en verheugd
met gastvrijheid als deugd.

“Wat ben je nog steeds klein,
maar oog je net zo fijn”,

zei tante met ontzag:
“als toen je zus net lag

in haar wiegje van hout.
Ze had het niet snel koud.

Ik haar dekens afsloeg,
al kwam ze veel te vroeg,

omdat ze het warm kreeg.
Waarop ik maar verzweeg

dat Emma ouder leek.
Haar huid was niet echt bleek,

maar zelfs schilferig droog.
Ze keek alert omhoog

als iemand binnenkwam,
in plaats van sloom of lam.

Ze had al flink wat haar
en leek al meer dan klaar
om voldragen te zijn.
Dat was natuurlijk fijn,

want we hadden zorgen
of de dag van morgen

wel kon toebehoren
aan zo vroeg geboren

broze zuigelingen.
Het zijn van die dingen

waar je nog steeds aan denkt.
Hoe God het leven schenkt.”

Zo verliepen maanden.
Ze zich veilig waanden

tot de priester langs kwam
en haar oom apart nam.

Emmy bleek verstoten,
waarop zij besloten

haar kalm te verzoeken
hen niet op te zoeken

nu ze vogelvrij was.
Met een milde grimas

toonde ze zich dankbaar
voor wat ze al voor haar

gastvrij hadden gedaan.
Ze wilde eervol gaan.

Het gevaar ingeschat,
koos ze het hazenpad,

schuilend in de bossen
om zich te verlossen

van een denkbaar onheil.
In haar nieuwe leefstijl

voelde ze zich bevreesd.
Ze was nooit vrij geweest.

Ze kreeg plots inzage
in een oude sage,

die ze als kind hoorde
en zich weer verwoordde.

Het ging over reuzen,
die leefden als geuzen

pal in Midden-Drenthe.
Deze corpulente

bliksemse giganten
belaagden passanten.

Het overkwam mensen
waar aan elkaar grenzen,

op de markegronden
alwaar schapen stonden,

zeven Drentse dorpen. Eenmaal onderworpen

werden rijken beroofd
en soldaten onthoofd.

Emmy vond het gezwam
toen ze de blaam vernam.

Ze had zelf geen bezit.
Ze zag de blanke pit

van de ruwe bolsters,
die met lege holsters

geharde vechtjassen
wisten te verrassen,

wat ze waarderen kon.
Vroom had Emmy Mammon

in schijn afgewezen
en dus niets te vrezen.

Mensen die rijk waren,
konden niet ervaren

enige sympathie
van de arme Emmy.

Als die werden beroofd
dan stootten ze hun hoofd

vanwege hun hebzucht;
een tekort aan godsvrucht.

Ze hoorde later pas
hoe het vervolgstuk was.

De wrede kolossen
hadden in de bossen

brutaal een meid ontvoerd
en daarna afgevoerd

rats naar hun plaggenhut.
Daar werd ze vol benut

en aldoor geslagen.
Ze moest daar behagen

de barse barbaren
zeven helse jaren.

Met deze legende
doorspekt met ellende

wilde de adelheid
de bestaande bangheid
voor het buitengebied
en wat daar soms geschiedt,

gewiekst aanwakkeren
om zo de zwakkeren

te demotiveren
de rug toe te keren

naar een strenge grondheer.
Het werkte elke keer

want Emmy school een tijd
uit angst voor haar vrijheid.

Dat deed ze eerst alleen.
Haar pa zei ooit gemeen,

toen Emmy bij een kroeg,
hem daar retorisch vroeg,

lopend in zijn kielzog:
“U beschermt me wel, toch?"

Haar vader zei toen laf:
“Dat hangt van de vraag af.

Kun je me bijbenen?”
Haar pijlers verdwenen.

Waren er nooit geweest.
Dat schokte haar het meest.

Haar existentiële
onzekerheid speelde.

Haar uitvlucht was de aard,
zei later Kierkegaard,

van de bange massa.
Schuilend in een dogma.

Werd door tegenslagen
meer grond weggeslagen?
